\documentclass{article}

\textwidth=6.5in
\textheight=8.5in
\voffset=-54pt
\oddsidemargin=0in
\evensidemargin=0in
\setlength{\parskip}{8pt}
\setlength{\parindent}{0pt}

\usepackage{amsmath, amssymb}
\usepackage{ifthen}

\usepackage[inline]{enumitem}
\setlist{topsep=0pt, parsep=8pt}

\usepackage[rgb,table]{xcolor}\convertcolorsUfalse
\usepackage{graphicx}

% change hyperref options
\PassOptionsToPackage{hyphens}{url}
\usepackage[bookmarks,colorlinks,linkcolor=black,citecolor=blue,pdfstartview=FitH,urlcolor=blue]{hyperref}
\hypersetup{pdfpagemode=UseNone}
\usepackage{bookmark}

\usepackage{graphicx}
\usepackage{multicol}
\usepackage{framed}
\usepackage{array}

\usepackage{icomma}

\usepackage{soul}

\usepackage{pgfplots}
\pgfplotsset{compat=1.18}  
\pgfplotscreateplotcyclelist{mycolor}{black, blue, red, green!70!black, magenta, purple, cyan, orange }  
\pgfplotsset{
  disabledatascaling,
  cycle list=tikzcolor, 
  axis lines=middle,
  every axis plot/.style={no marks, thick},
  every axis/.style={
    enlarge x limits=false,
    enlarge y limits=auto,
    xlabel style={at={(current axis.right of origin)},anchor=west},
    ylabel style={at={(current axis.above origin)},anchor=south},
    % extra x and y ticks for asymptotes
    extra x tick style={grid=major, grid style={dashed,black}},
    extra y tick style={grid=major, grid style={dashed,black}},
    /pgf/number format/1000 sep={},
  },    
  tick label style = {font=\footnotesize, fill=\currentbgcolor, inner sep=0pt},    
  xticklabel shift=1pt,    
  scaled ticks=false,
  scaled y ticks=false,
  unbounded coords=jump,
  samples=101,
  minor tick length=0,
  scatter/classes={
    open={mark=*,draw=black,fill=white, mark size=2, thin},
    closed={mark=*,draw=black,fill=black, mark size=2, thin}
  },
  width=3in,
}
\pgfplotsset{open/.style={only marks, mark=*, fill=white, thin, mark size=2}}
\pgfplotsset{closed/.style={only marks, mark=*, draw=black, fill=black,
    thin, mark size=2}}
\pgfplotsset{labeled/.style={nodes near coords, only marks, mark=*, draw=black, fill=black, thin, mark size=2, point meta=explicit symbolic}}

\usepackage{tikz}
\usetikzlibrary{
  matrix,%
  % enables the snake paths
  decorations.pathmorphing,%
  % enables bigger arrows
  decorations.markings,%
  decorations.pathreplacing,%
  decorations.text,%
  calc,%
  patterns,%
  shapes.geometric,%
  arrows,
  decorations.shapes,
  positioning,
  plotmarks,
  intersections
}
\tikzset{>=latex} % sets default arrow type in tikz
\tikzset{font=\normalsize}
\tikzset{mark size=1.5pt, mark options=thin}
\tikzset{pin distance=4pt,
  every pin edge/.style={<-, thin, shorten <= -2pt}}

\interfootnotelinepenalty=10000
\renewcommand{\thefootnote}{(\arabic{footnote})}

\usepackage{multirow, bigdelim}

% change to \usepackage[sols]{handout} to see solutions
\usepackage[sols, grading, workshop]{handout}

\newcommand\iscurrentenumi[1]{%
  \ifthenelse{\equal{#1}{\theenumi}}{}{#1}}
\setenumerate[1]{ref=\#\arabic*}
\setenumerate[2]{ref=\protect\iscurrentenumi{\theenumi}(\alph*)}
\newcommand\enumiiprefix[2]{%
  \ifthenelse{{\equal{#1}{\theenumi}}\AND{\equal{#2}{(\alph{enumii})}}}{}{}%
  \ifthenelse{{\equal{#1}{\theenumi}}\AND{\NOT{\equal{#2}{(\alph{enumii})}}}}{#2}{}%
  \ifthenelse{\equal{#1}{\theenumi}}{}{#1#2}%
}
\setenumerate[3]{ref=\protect\enumiiprefix{\theenumi}{(\alph{enumii})}\roman*}

\makeatletter

% underline links               
\let\@href=\href
\renewcommand{\href}[2]{\@href{#1}{\ul{#2}}}

\makeatother

\definecolor{purple}{rgb}{0.5,0,1.0}

\newcounter{imagecounter}
\newcounter{columncounter}
\newenvironment{imagetable}[3][]{%
  \setcounter{imagecounter}{0}
  \setcounter{columncounter}{#2}
  \addtocounter{columncounter}{-1}
  \newcolumntype{i}{>{\raisebox{-1ex-\height}{\makebox[1em]{\hfill\stepcounter{imagecounter}(#3{imagecounter})}}\hspace{4pt}\begin{minipage}[t]{\linewidth/#2-1em-\tabcolsep*3}\arraybackslash\vspace{0pt}}l<{\end{minipage}}}
\begin{tabular}{i*{\value{columncounter}}{#1i}}}
{\end{tabular}}  

\newcommand{\doedfinity}[2][\newpage]{
  \begin{problem}[#1]
    Do the Edfinity assignment ``Problem Set #2''.

    We encourage you to write down any scratch work here, as well as
    things you want your future self to remember. Nothing you write
    here will be graded; it's just for you!
  \end{problem}
}

\newcommand{\psrnewpage}{\newpage}
\newcommand{\psreflection}[2][]{

  \ifthenelse{\not\equal{#1}{nocite}}{
    \begin{enumerate}[resume]
    \item
      \begin{problem}[1in]
        Please cite any help you got on this problem set:
        \begin{itemize}[nosep]
        \item If you worked with other students, please list their names here.
        \item If you used resources other than those on our Canvas site, please list them here.
        \end{itemize}
      \end{problem}

    \end{enumerate}
  }

  \probonly{%
    #2

    \psrnewpage

    \emph{Reflection space.} It's a good habit to take a few minutes at the end of each pset to reflect on what you learned and jot down notes to your future self. Here are a few reflection questions to get you started. Nothing you write here will be graded; it's just for you!
    \begin{itemize}
    \item Take a few minutes to look back at the problems. How could you explain to someone else how to do each problem? What was your strategy, and how did you know to use that strategy? If there are problems where you don't feel able to answer this, write down which ones they are. Come back to them in a few days when we post the homework solutions!

      \vspace{1.5in}

    \item Take a look back at the learning objectives on today's worksheet; how are you feeling about each one? If there are ones you know you need more practice with, write those down here.

      \vspace{1.5in}

    \item What did you find most challenging about this material? Do you have any lingering questions about it? If so, write them down here so you don't forget them. At some point, stop by office hours or MQC to ask!

      \vfill
    \end{itemize}      
  }
}

\usepackage{fancyhdr}
\lhead{\textcolor{white}{This content is protected and may not be shared, uploaded, or distributed.}}
\rhead{}
\rfoot{\vspace*{\fill}\footnotesize\textcopyright{} \emph{Harvard Math Ma}}
\renewcommand{\headrulewidth}{0pt}
\pagestyle{fancy}
\begin{document}

\begin{center}
  \large\textsc{Problem Set 16 - Continuity}
\end{center}

\begin{enumerate}
\item  \note{
    \begin{itemize}[nosep]
    \item Recognize removable discontinuities, jumps, and vertical asymptotes graphically.
    \item On a graph, identify points where $f$ is not continuous / not differentiable.
    \item Find discontinuities of a rational function.
    \item Make a sign chart for something like $f'(x) = \dfrac{x - 5}{x + 8}$, and determine where $f$ is increasing.
    \end{itemize}
}

  \doedfinity{16}

\item  \begin{grading}
    1 point per part
  \end{grading}

  Suppose $f$ is a function defined on $(-\infty, \infty)$, and you're given the following information about $f$.
  \begin{alignat*}{3}
    \lim_{x \to 2^-} f(x) &= 3 \hspace{1in} &
    \lim_{x \to 4^-} f(x) &= 1 \hspace{1in} &
    \lim_{x \to 5^-} f(x) &= -1 \\
    \lim_{x \to 2^+} f(x) &= 3 & \lim_{x \to 4^+} f(x) &= 5 & \lim_{x \to
      5^+} f(x) &= -1  \\
    f(2) &= 4 
    & f(4) &= 1 
    & f(5) &= -1 
  \end{alignat*}
  Decide whether $f(x)$ is continuous at each of the following values of $x$; if not, what kind of discontinuity (removable discontinuity, jump discontinuity, or vertical asymptote) does $f$ have there? Explain briefly.

  \begin{enumerate}
  \item
    \begin{problem}[0.5in]
      $x = 2$
    \end{problem}

    \begin{solution} % Jason Bland
      $\displaystyle \lim_{x \to 2} f(x)$ exists (it equals 3), but it doesn't equal $f(2)$, so $f$ has a \fbox{removable discontinuity} at $x = 2$.
    \end{solution}

  \item
    \begin{problem}[0.5in]
      $x = 4$
    \end{problem}

    \begin{solution} % Jason Bland
      $\displaystyle \lim_{x \to 4^-} f(x)$ and $\displaystyle \lim_{x \to 4^+} f(x)$ exist but aren't equal to each other, so $f$ has a \fbox{jump discontinuity} at $x = 4$.
    \end{solution}

  \item
    \begin{problem}[0.5in]
      $x = 5$
    \end{problem}

    \begin{solution} % Jason Bland
      $\displaystyle \lim_{x \to 5} f(x) = -1$, which is equal to $f(5)$, so $f$ is \fbox{continuous} at $x = 5$.
    \end{solution}

  \end{enumerate}  

\item  \begin{grading}
    1 point each
  \end{grading}

  \begin{enumerate}
  \item

    \begin{problem}[\vfill]
      Sketch an example of a function $f$ that's continuous at $x = 3$ but not differentiable at $x = 3$, or say that it's impossible. 
    \end{problem}

    \begin{solution}
    
      \begin{tikzpicture}
        \begin{axis}[xtick={3}, ytick=\empty, width=2in]
          \addplot[domain=1:5] {-abs(x-3) + 3};
        \end{axis}
      \end{tikzpicture}
    \end{solution}

  \item
    \begin{problem}[\vfill\newpage]
      Sketch an example of a function $f$ that's differentiable at $x = 3$ but not continuous at $x = 3$, or say that it's impossible. 
    \end{problem}

    \begin{solution}
      This is impossible.
    \end{solution}

  \item
    \begin{problem}[\vfill]
      Sketch an example of a function $f$ that's both continuous at $x = 3$ and differentiable at $x = 3$, or say that it's impossible. 
    \end{problem}

    \begin{solution}
    
      \begin{tikzpicture}
        \begin{axis}[xtick={3}, ytick=\empty, width=2in]
          \addplot[domain=1:5] {-(x-3)^2 + 3};
        \end{axis}
      \end{tikzpicture}
    \end{solution}

  \end{enumerate}

\item  \begin{grading}
    3 points: 1 point for putting both $4$ and $-7$ on their sign chart, 1.5 for testing the regions correctly, 0.5 for having a conclusion consistent with their sign chart.

    If they don't realize $f$ could change sign at $-7$, they should still realize $f$ is increasing on $(4, \infty)$; give them 1.5 points for this. 
  \end{grading}

  \begin{problem}[\nstretch{2}]
    Suppose $f$ has domain $(-\infty, \infty)$ and $f'(x) = \dfrac{(x - 4)(x^2 + 1)}{x + 7}$. Make a sign chart for $f'$, and use this to determine where $f$ is increasing and where $f$ is decreasing.
  \end{problem}

  \begin{solution}
    First, we need to determine the values for which $f'$ is zero or undefined.
    Looking at the denominator, $f'(x)$ is undefined when $x=-7$. Looking at the numerator,
    $f'(x)$ is zero when $x=4$. These values go on our sign chart.
      \begin{center}
        \begin{tikzpicture}
          \draw[<->] (-3, 0) -- (3, 0);
          \node at (-3.5, 0) {$f'(x)$};
          \draw[shift={(-1,0)},color=black] (0pt,3pt) -- (0pt,-3pt) node[below] {$-7$};
          \draw[shift={(1,0)},color=black] (0pt,3pt) -- (0pt,-3pt) node[below] {$4$};
          \node[circle, fill=black, inner sep=2pt] at (1,0) {};
          \node[circle, fill=white, draw=black, inner sep=2pt] at (-1,0) {};
        \end{tikzpicture}
      \end{center}
      To determine the sign of $f'(x)$, we test one $x$-value in each interval.
      \begin{alignat*}{3}
      f'(-8) &= \frac{((-8)-4)((-8)^2+1)}{(-8)+7} &= 
      \frac{(\text{negative})(\text{positive})}{\text{negative}} &= \text{positive} \\
      f'(0) &= \frac{((0)-4)((0)^2+1)}{(0)+7} &= 
      \frac{(\text{negative})(\text{positive})}{\text{positive}} &= \text{negative} \\
      f'(5) &= \frac{((5)-4)((5)^2+1)}{(5)+7} &= 
      \frac{(\text{positive})(\text{positive})}{\text{positive}} &= \text{positive} \\
      \end{alignat*}
      We can then add this information to our sign chart.
      \begin{center}
        \begin{tikzpicture}
          \draw[<->] (-3, 0) -- (3, 0);
          \node at (-3.5, 0) {$f'(x)$};
          \draw[shift={(-1,0)},color=black] (0pt,3pt) -- (0pt,-3pt) node[below] {$-7$};
          \draw[shift={(1,0)},color=black] (0pt,3pt) -- (0pt,-3pt) node[below] {$4$};
          \node[circle, fill=black, inner sep=2pt] at (1,0) {};
          \node[circle, fill=white, draw=black, inner sep=2pt] at (-1,0) {};
          \node[above] at (-2,0) {$+$};
          \node[above] at (0,0) {$-$};
          \node[above] at (2,0) {$+$};
        \end{tikzpicture}
      \end{center}

      Thus we can determine that $f$ is increasing on $(-\infty, -7)$ and $(4, \infty)$,
      since $f'$ is positive on those intervals, and $f$ is decreasing on $(-7, 4)$, since
      $f'$ is negative on that interval.
      
  \end{solution}

\item \label{discontinuities-and-horizontal-asymptotes-1}

  Let $f(x) = \dfrac{3x^2 + 9x}{x^2 - 9}$.

  \begin{enumerate}
  \item \label{discontinuities-and-horizontal-asymptotes-1-discontinuities}

    \begin{grading}
      1 point
    \end{grading}

    \begin{problem}[0.75in]
      Find all values in $(-\infty, \infty)$ at which $f$ is discontinuous. 
    \end{problem}

    \begin{solution}
      Since $f$ is defined as a single algebraic formula, $f$ is continuous
      as long as $f(x)$ is defined. First, let's factor the numerator and denominator.
      \[f(x) = \frac{3x(x+3)}{(x+3)(x-3)}.\]
      Looking at the denominator, the only discontinuities
      are at $x = -3$ and $x=3$. 
    \end{solution}

  \item

    \begin{grading}
      5.5 points:
      \begin{itemize}[nosep]
      \item 1.5 for $\displaystyle \lim_{x \to -3} f(x)$
      \item 1.5 for $\displaystyle \lim_{x \to 3^-} f(x)$
      \item 1.5 for $\displaystyle \lim_{x \to 3^+} f(x)$
      \item 1 for classifying the discontinuities 
      \end{itemize}
    \end{grading}

    \begin{problem}[\vfill\newpage]
      Classify each discontinuity (that is, say whether it's a removable discontinuity, jump discontinuity, or vertical asymptote) using limits. If a limit does not exist, please calculate both one-sided limits. Be sure to show any work you use to compute the relevant limits, and organize your work so that your reader can follow it easily. 
    \end{problem}

    \begin{solution}
      Let's start with $x=-3$. If we try to evaluate $\lim\limits_{x \to -3} f(x)$ using
      direct substitution, we get $\frac{0}{0}$, so we should try to factor and cancel.
      \begin{align*}
        \lim_{x \to -3} f(x) &= \lim_{x \to -3} \frac{3x(x+3)}{(x+3)(x-3)} \\
        &= \lim_{x \to -3} \frac{3x}{x-3} \\
        &= \frac{3(-3)}{(-3)-3} \\
        &= \frac{3}{2}.
      \end{align*}
      Since $\lim\limits_{x \to -3} f(x)$ exists but is not equal to $f(-3)$, this
      is a removable discontinuity.

      Now let's look at $x=3$. Whenever $x \neq -3$, $f(x) = \frac{3x}{x-3}$.
      If we attempt to evaluate $\lim\limits_{x \to 3} f(x)$ using
      direct substitution, we get $\frac{9}{0}$, which indicates that we are likely
      dealing with a vertical asymptote. Let's look at the one-sided limits.
      As $x$ approaches $3$ from the left, the numerator is positive and the denominator
      is negative, so $\lim\limits_{x \to 3} f(x) = -\infty$. As $x$ approaches $3$ from
      the right, the numerator is positive and the denominator is positive, so 
      $\lim\limits_{x \to 3} f(x) = \infty$. Since at least one of the one-sided limits
      is $\pm \infty$, this is an infinite discontinuity aka a vertical asymptote.
    \end{solution}

  \item \label{discontinuities-and-horizontal-asymptotes-1-HA}

    \begin{grading}
      2.5 points: 2 for calculating $\displaystyle \lim_{x \to \infty} f(x)$ and $\displaystyle \lim_{x \to -\infty} f(x)$ correctly and 0.5 for stating what the horizontal asymptote is
    \end{grading}

    \begin{problem}[\vfill]
      Find all horizontal asymptotes of $f$ (if it has any). Be sure to show your calculation of the relevant limits. 
    \end{problem}

    \begin{solution}
      To find the horizontal asymptote(s) of $f$, we need to evaluate the limits
      as $x$ approaches $\pm \infty$.

      First let's consider the limit as $x$ goes to $-\infty$.
      \begin{align*}
        \lim_{x \to -\infty} f(x) &= \lim_{x \to -\infty} \frac{3x}{x-3} \\
        \intertext{The dominant term in the denominator as $x \to -\infty$ is $x$.}
        &= \lim_{x \to -\infty} \frac{3x}{x} \\
        &= \lim_{x \to -\infty} 3 \\
        &= 3.
      \end{align*}

      Now let's consider the limit as $x$ goes to $\infty$. 
      \begin{align*}
        \lim_{x \to \infty} f(x) &= \lim_{x \to \infty} \frac{3x}{x-3} \\
        \intertext{The dominant term in the denominator as $x \to \infty$ is $x$.}
        &= \lim_{x \to \infty} \frac{3x}{x} \\
        &= \lim_{x \to \infty} 3 \\
        &= 3.
      \end{align*}
      This shows that the graph of $y=f(x)$ has one horizontal asymptote, $y=3$.
    \end{solution}

  \item 

    \begin{grading}
      2 points: 1 for the choice, 1 for the explanation 
    \end{grading}

    \begin{problem}[\nstretch{0.35}]
      Using just what you found out in \ref{discontinuities-and-horizontal-asymptotes-1-discontinuities} - \ref{discontinuities-and-horizontal-asymptotes-1-HA}, decide which of the following shows the graph of $f$. Explain briefly.

      \begin{imagetable}{3}{\Alph}
        \begin{tikzpicture}
          \begin{axis}[width=2.25in, xtick={-3, 3}, ymin=-12, ymax=12]
            \addplot+[domain=-6:6] {2+6/((x-1)^2+1)};%
            \addplot+[open] coordinates {(-3, 40/17) (3, 3.2)};% 
          \end{axis}
        \end{tikzpicture} &
        \begin{tikzpicture}
          \begin{axis}[width=2.25in, xtick={-3, 3}, ymin=-12, ymax=12]
            \addplot+[domain=-6:2.9]{-4/3*x - 3/(x-3)-3};% 
            \addplot+[domain=3.1:6]{-x - 1.5/(x-3) + 1.25};%
            \addplot+[open] coordinates { (-3, 1.5) };% 
          \end{axis}
        \end{tikzpicture} &
        \begin{tikzpicture}
          \begin{axis}[width=2.25in, xtick={-3, 3}, ymin=-12, ymax=12]
            \addplot+[domain=-6:2.9]{3 + 9/(x-3)};% 
            \addplot+[domain=3.1:6]{3 + 9/(x-3)};%
            \addplot+[open] coordinates { (-3, 1.5) };% 
          \end{axis}
        \end{tikzpicture}
        \vspace{12pt} \\
        \begin{tikzpicture}
          \begin{axis}[width=2.25in, xtick={-3, 3}, ymin=-12, ymax=12]
            \addplot+[domain=-6:-3.1]{3 - 9/(x+3)};% 
            \addplot+[domain=-2.9:6]{3 - 9/(x+3)};%
            \addplot+[open] coordinates { (3, 1.5) };% 
          \end{axis}
        \end{tikzpicture} &
        \begin{tikzpicture}
          \begin{axis}[width=2.25in, xtick={-3, 3}, ymin=-12, ymax=12]
            \addplot+[domain=-6:-3.05]{2 + 1/(x+3)+1/(x-3)};% 
            \addplot+[domain=-2.95:2.95]{2 + 1/(x+3)+1/(x-3)};%
            \addplot+[domain=3.05:6]{2 + 1/(x+3)+1/(x-3)};%
          \end{axis}
        \end{tikzpicture} &
        \begin{tikzpicture}
          \begin{axis}[width=2.25in, xtick={-3, 3}, ymin=-12, ymax=12]
            \addplot+[domain=-6:-3.05]{x + 1/(x+3)^2+1/(x-3)^2};% 
            \addplot+[domain=-2.95:2.95]{x + 1/(x+3)^2+1/(x-3)^2};% 
            \addplot+[domain=3.05:6]{x + 1/(x+3)^2+1/(x-3)^2};% 
          \end{axis}
        \end{tikzpicture} 
      \end{imagetable}
    \end{problem}

    \begin{solution}
      We know that there is a hole at $x=-3$, a vertical asymptote at $x=3$, and
      a horizontal asymptote at $y=3$. This matches graph (C).
    \end{solution}
  \end{enumerate}

\item  \begin{grading} 
    5 points:
    \begin{itemize}[nosep]
    \item 4 for correctly simplifying the difference quotient $\dfrac{ f(x+h) - f(x) }{h}$ (give partial credit if the simplification is partially correct)
    \item 0.5 for using limit notation correctly (writing $\displaystyle \lim_{h \to 0}$ on each line until they actually take the limit)
    \item 0.5 for correct limit
    \end{itemize}
  \end{grading} 

  \probonly{\emph{Next class, the definition of derivative will be important again, so this is a reminder.}}

  \begin{problem}[\stretch{2}]
    Let $f(x) = \dfrac{1}{x^2}$. Use the definition of the derivative to calculate $f'(x)$.
  \end{problem}

  \begin{solution}  
    \begin{align*} 
      f'(x ) &= \displaystyle \lim_{h \rightarrow 0} \frac{ f(x+h) - f(x) }{x+h-x}\\
      & = \displaystyle \lim_{h \rightarrow 0} \frac{ \frac{1}{(x+h)^2}-\frac{1}{x^2}} {h} \\
      &= \displaystyle \lim_{h \rightarrow 0} \frac{ \frac{ x^2 - (x+h)^2}{x^2(x+h)^2}}{h} \\
      &= \displaystyle \lim_{h \rightarrow 0} \frac{ x^2 -(x^2 +2xh+h^2)}{x^2(x+h)^2} \cdot \frac{ 1}{h} \\
      &= \displaystyle \lim_{h \rightarrow 0} \frac{- 2xh -h^2}{x^2(x+h)^2} \cdot \frac{1}{h} \\
      &= \displaystyle \lim_{h \rightarrow 0}\frac{ h(-2x-h)}{x^2(x+h)^2}\cdot \frac{1}{h} \\
      &= \displaystyle \lim_{h \rightarrow 0} \frac{ -2x-h}{x^2(x+h)^2}\\
      &=\frac{-2x}{x^4} \\
      &=\boxed{ \frac{-2}{x^3}}
    \end{align*} 
  \end{solution}  

\end{enumerate}
\psreflection{For next class, read \href{https://people.math.harvard.edu/~jjchen/mathma/docs/handouts/Hughes-Hallet-3-1.pdf}{\S 3.1 of \emph{Calculus} by Hughes-Hallett et. al} through Example 4.}
\end{document}
